\section{Methodology}
\label{sec:methodology}

The framework proceeds through five sequential stages:
candidate generation, spectral clustering, MCA-based demand weighting,
proportional budget allocation, and clustered \mbox{p-median} optimization.
Each stage is described below, together with the design rationale
that distinguishes each algorithmic choice from simpler alternatives.

% ---------------------------------------------------------------
\subsection{Pedestrian Network and Candidate Generation}
\label{ssec:network}

All distance computations are performed on the pedestrian street network
downloaded from OpenStreetMap~\cite{osm2025} and processed with
OSMnx~\cite{boeing2017osmnx}.
Census-block centroids and existing store locations are snapped to their
nearest graph node before shortest-path computation.
Let $G = (V, E)$ denote the projected walking graph and let $d_{u,v}$
be the shortest-path distance between nodes $u, v \in V$.

The first screening stage identifies blocks that are potentially underserved
by the existing store network.
For each census block~$i$, let $\delta_i$ denote its shortest-path distance
to the nearest existing store.
A block is classified as a \emph{candidate} if
\begin{equation}
  \delta_i > \tau \cdot \bar{d}_{\text{nn}},
  \label{eq:threshold}
\end{equation}
where $\bar{d}_{\text{nn}}$ is the empirical average nearest-neighbor distance
among all existing stores and $\tau$ is a multiplier calibrated from data.
Setting $\tau = 1.0$—so that the threshold matches the observed inter-store
spacing—produces a candidate set whose size is directly tied to the
existing coverage density rather than to an arbitrary absolute cutoff.
Sensitivity of coverage rate to $\tau \in \{0.5, 0.75, 1.0, 1.25, 1.5\}$
is reported in \cref{sec:sensitivity}.

% ---------------------------------------------------------------
\subsection{Graph-Based Spectral Clustering}
\label{ssec:spectral}

Computing all-pairs pedestrian distances is impractical for large instances.
Instead, a sparse affinity matrix is constructed in two steps.
For each census block, the $K = 12$ nearest neighbors in Euclidean space are
identified; for those candidate pairs only, the actual road-network distance
is computed.
The affinity between blocks $i$ and $j$ is defined by a Gaussian kernel:
\begin{equation}
  w_{i,j} = \exp\!\left(-\frac{d_{i,j}^{2}}{2\sigma^{2}}\right),
  \label{eq:affinity}
\end{equation}
where $\sigma$ is the median observed road-network distance among all evaluated
neighbor pairs.
The resulting sparse matrix $W$ defines a weighted similarity graph
whose edges reflect actual network connectivity rather than geometric proximity.

The normalized graph Laplacian is:
\begin{equation}
  L = D^{-1/2}(D - W)D^{-1/2},
  \label{eq:laplacian}
\end{equation}
where $D$ is the diagonal degree matrix of $W$.
The initial cluster count $k^*$ is selected by the eigengap heuristic:
\begin{equation}
  k^* = \arg\max_k \left(\lambda_{k+1} - \lambda_k\right),
  \label{eq:eigengap}
\end{equation}
using the smallest eigenvalues of $L$~\cite{vonluxburg2007tutorial}.
\Cref{fig:eigengap} shows the eigenvalue spectrum and eigengap profile
for the student-area instance ($k^* = 16$); the result is confirmed by a
secondary silhouette analysis (\cref{fig:silhouette}).

The raw partition typically contains several very small territories.
To produce operationally deployable regions, clusters with fewer than
200 census blocks are merged into the nearest major cluster using
road-network distance whenever available and Euclidean distance as a fallback.
\Cref{fig:clusters} compares the raw $k^*$-cluster partition against
the final merged solution for the student-area instance.

The $K = 12$ neighbor count was validated by comparing eigengap stability
across $K \in \{6, 8, 10, 12, 15, 18\}$: lower values produce less stable
eigengaps and more fragmented clusters, while the signal is consistent
for $K \geq 12$ (\cref{fig:k_sensitivity}).

\begin{figure}[t]
  \centering
  \begin{subfigure}[b]{0.49\textwidth}
    \includegraphics[width=\textwidth]{spectral_eigengap_decision}
    \caption{Eigenvalue spectrum and eigengap ($k^*=16$, student area).}
    \label{fig:eigengap}
  \end{subfigure}\hfill
  \begin{subfigure}[b]{0.49\textwidth}
    \includegraphics[width=\textwidth]{spectral_silhouette_support}
    \caption{Silhouette coefficient for $k$ values around $k^*$.}
    \label{fig:silhouette}
  \end{subfigure}
  \caption{Spectral cluster-count selection for the student-area instance.
           The eigengap selects $k^*=16$; silhouette analysis provides
           secondary confirmation.}
  \label{fig:cluster_selection}
\end{figure}

\begin{figure}[t]
  \centering
  \includegraphics[width=\textwidth]{graph_spectral_clusters_before_after_merge}
  \caption{Raw 16-cluster partition (left) and the six-cluster solution
           after the small-cluster merge rule (right) for the student area.
           Road-network affinity preserves urban connectivity; the merge rule
           converts fragmented micro-territories into deployment-ready zones.}
  \label{fig:clusters}
\end{figure}

% ---------------------------------------------------------------
\subsection{Socio-Urban Demand Weighting via MCA}
\label{ssec:mca}

Population count alone is an insufficient demand proxy for retail location
because blocks of equal population can differ sharply in infrastructure quality,
transport access, and economic activity.
MCA is used instead of PCA because the source variables are categorical
or categorized~\cite{greenacre2017correspondence,greenacre2006multiple}.
Each continuous block-level variable is discretized into quintiles when
it has sufficient cardinality; otherwise, original categories are retained.
Missing values are explicitly encoded as an additional category so that
no observation is discarded.
The final indicator matrix contains 102 dummy columns derived from
33 source variables covering age structure, disability, schooling,
housing conditions, household services, and urban infrastructure
(street signage, bicycle lanes, pedestrian crossings, ramps,
public lighting, collective transport, and trees).

The resulting MCA coordinates are denoted $\text{MCA}_1, \text{MCA}_2, \ldots$
After Greenacre's correction~\cite{greenacre1991interpreting},
$\text{MCA}_1$ explains approximately 72\% of adjusted inertia
in the student-area instance (\cref{fig:mca_inertia}).
Infrastructure and urban-environment indicators---street signage, bicycle
infrastructure, pedestrian crossings, ramps, transport access, education,
and housing services---load most strongly on this dimension.
Categories reflecting higher schooling and more complete urban infrastructure
appear on the positive end of $\text{MCA}_1$, making it interpretable as a
socio-urban quality gradient rather than a purely demographic axis.

To prevent arbitrary sign from entering the optimization weights,
$\text{MCA}_1$ is oriented so that higher schooling corresponds to
positive values; if the raw extraction is inverted, all scores are
multiplied by $-1$.
The normalized demand modifier is then:
\begin{equation}
  \tilde{m}_i =
    \frac{\text{MCA}_{1,i} - \min(\text{MCA}_1)}
         {\max(\text{MCA}_1) - \min(\text{MCA}_1)},
  \label{eq:mca_norm}
\end{equation}
which rescales the gradient to $[0, 1]$ without altering relative differences.
\Cref{fig:mca_maps} confirms that $\text{MCA}_1$ scores exhibit strong
spatial continuity, which validates their use as a geographically structured
demand modifier.
\Cref{fig:mca_by_cluster} shows that spectral clusters differ not only
in geometry but also in their MCA profiles, reinforcing the rationale
for solving the optimization model within each territory independently.

\begin{figure}[t]
  \centering
  \begin{subfigure}[b]{0.49\textwidth}
    \includegraphics[width=\textwidth]{mca_explained_components}
    \caption{Raw and Greenacre-corrected explained inertia by component.}
    \label{fig:mca_inertia}
  \end{subfigure}\hfill
  \begin{subfigure}[b]{0.49\textwidth}
    \includegraphics[width=\textwidth]{mca_dims_by_cluster_boxplot}
    \caption{MCA\textsubscript{1,2,3} by spatial cluster.}
    \label{fig:mca_by_cluster}
  \end{subfigure}
  \caption{MCA component structure and its relationship to the spectral
           territory partition. The first corrected dimension concentrates
           most of the socio-urban variation; clusters differ in MCA profile
           as well as in spatial extent.}
  \label{fig:mca_structure}
\end{figure}

\begin{figure}[t]
  \centering
  \includegraphics[width=\textwidth]{mca123_spatial_maps}
  \caption{Choropleth maps of $\text{MCA}_1$ (left), $\text{MCA}_2$ (center),
           and $\text{MCA}_3$ (right) for the student-area instance.
           The strong spatial continuity of $\text{MCA}_1$ validates its
           role as a geographically structured demand weight.}
  \label{fig:mca_maps}
\end{figure}

% ---------------------------------------------------------------
\subsection{Demand-Proportional Budget Allocation}
\label{ssec:allocation}

The original formulation allocated a fixed number of new openings ($p_{\text{new}}$)
to each cluster regardless of its size or demand.
In a six-cluster partition where cluster sizes range from 212 to 1,448 blocks
(\cref{tab:allocation}), a fixed allocation of seven openings per cluster
implies radically different per-demand coverage ratios across territories.
The largest cluster receives one opening per 184 blocks; the smallest,
one per 30.

To correct this, a proportional allocation rule is introduced:
\begin{equation}
  p_c = \max\!\left(1,\; \left\lfloor
    P_{\text{total}} \cdot \frac{\sum_{i \in C_c} w_i}
                                {\sum_{i \in I} w_i}
  \right\rceil\right),
  \label{eq:proportional}
\end{equation}
where $C_c$ is the set of demand blocks in cluster~$c$, $w_i$ is the
demand weight defined in \cref{eq:demand_weight} below,
and the allocation is rounded to the nearest integer with a minimum of one.
Integer remainder is distributed to the largest cluster.
\Cref{fig:allocation} shows the shift from fixed to proportional allocation
for all three study areas; the largest clusters receive substantially more
openings, which translates into measurable coverage gains in
blocks-per-opening terms.

\begin{table}[t]
  \centering
  \caption{Demand-proportional budget allocation versus the original
           fixed-per-cluster rule (student-area instance).}
  \label{tab:allocation}
  \begin{tabular}{rrrrrrr}
\toprule
\textbf{Cluster}
  & \textbf{Demand nodes}
  & \textbf{Population}
  & \textbf{Demand share (\%)}
  & \textbf{Fixed alloc.}
  & \textbf{Prop.\ alloc.} \\
\midrule
0 & 1{,}282 & 112{,}486 & 27.4 & 7 & 11 \\
1 & 1{,}182 & 102{,}309 & 25.1 & 7 & 11 \\
2 & 1{,}107 &  95{,}921 & 23.2 & 7 & 10 \\
3 &   590 &  49{,}175 & 11.9 & 7 &  5 \\
4 &   358 &  30{,}732 &  7.5 & 7 &  3 \\
5 &   212 &  19{,}836 &  4.9 & 7 &  2 \\
\midrule
\textbf{Total} & \textbf{4{,}731} & \textbf{410{,}459} & \textbf{100.0}
  & \textbf{42} & \textbf{42} \\
\bottomrule
\end{tabular}

\end{table}

\begin{figure}[t]
  \centering
  \includegraphics[width=0.7\textwidth]{fixed_vs_proportional_allocation}
  \caption{Fixed (original) versus proportional (revised) opening allocation
           by cluster for the student-area instance. The proportional rule
           concentrates the budget where demand is heaviest.}
  \label{fig:allocation}
\end{figure}

% ---------------------------------------------------------------
\subsection{Clustered P-Median Formulation}
\label{ssec:pmedian}

For each cluster $c$ the model is solved independently.
Let $I_c$ and $J_c$ denote the demand blocks and candidate sites
within cluster~$c$, let $E_c$ be the existing stores whose service
radii overlap with~$I_c$, and let $P_c = E_c \cup J_c$.
For each demand block~$i$, let $P_i = \{p \in P_c : d_{i,p} \leq D_{\max}\}$
be the set of feasible service points.
Let $\mathcal{C}_c \subset J_c \times J_c$ be the set of candidate pairs
separated by less than $S_{\min}$ (the minimum separation constraint
that prevents demand cannibalization).

The demand weight for block~$i$ is:
\begin{equation}
  w_i = p_i (1 + \beta \tilde{m}_i),
  \label{eq:demand_weight}
\end{equation}
where $p_i$ is the population of block~$i$ and $\beta \geq 0$ controls
the influence of the MCA modifier.
Setting $\beta = 0$ reduces the model to pure population weighting;
sensitivity to $\beta \in \{0.0, 0.1, 0.25, 0.5, 1.0\}$ is analyzed in
\cref{sec:sensitivity}.

Binary decision variables are:
$x_{i,p} = 1$ if block~$i$ is assigned to service point~$p$;
$y_j = 1$ if candidate~$j$ is opened; and
$u_i = 1$ if block~$i$ remains uncovered.
The objective minimizes total weighted assignment distance plus a penalty
for uncovered demand:
\begin{equation}
  \min \sum_{i \in I_c} \sum_{p \in P_i} w_i d_{i,p} x_{i,p}
       + \lambda \sum_{i \in I_c} w_i u_i.
  \label{eq:objective}
\end{equation}
Subject to:
\begin{align}
  \sum_{p \in P_i} x_{i,p} + u_i &= 1, && \forall i \in I_c,
  \label{eq:assignment} \\
  \sum_{j \in J_c} y_j &= p_c,
  \label{eq:budget} \\
  y_j + y_k &\leq 1, && \forall (j, k) \in \mathcal{C}_c,
  \label{eq:separation} \\
  x_{i,j} &\leq y_j, && \forall i \in I_c,\; j \in J_c \cap P_i,
  \label{eq:implication}
\end{align}
with $x_{i,p}, y_j, u_i \in \{0, 1\}$.
All subproblems are solved to optimality in the reported experiments
(\cref{sec:results}).

% ---------------------------------------------------------------
\subsection{Parameter Calibration}
\label{ssec:calibration}

A key weakness of the original submission was that parameters
$D_{\max}$, $S_{\min}$, $\beta$, and $\tau$ were fixed without empirical
justification.
This paper calibrates them jointly through a Pareto-frontier grid search
over 81 combinations of $D_{\max} \in \{300, 366, 450\}$~m,
$S_{\min} \in \{200, 240, 300\}$~m, and $p_c \in \{5, 7, 9\}$
per cluster (for three budget levels), repeated for
$\beta \in \{0.1, 0.25, 0.5\}$.

The Pareto criterion jointly optimizes coverage rate and weighted mean
assignment distance; parameter combinations on the frontier
are identified as calibration candidates.
The chosen parameters are selected from the frontier as those maximizing
coverage without simultaneously increasing weighted distance
beyond the frontier's convex hull (\cref{fig:pareto}).
\Cref{tab:parameters} reports the calibrated values for each study area.

\begin{table}[t]
  \centering
  \caption{Calibrated model parameters by study area.
           All values were selected from the Pareto grid rather than
           fixed \textit{a priori}.}
  \label{tab:parameters}
  \begin{tabular}{lllll}
\toprule
\textbf{Parameter} & \textbf{Meaning}
  & \textbf{Student} & \textbf{City} & \textbf{Industrial} \\
\midrule
$D_{\max}$ (m)
  & Service radius
  & 300 & 366 & 450 \\
$S_{\min}$ (m)
  & Min.\ candidate separation
  & 200 & 200 & 200 \\
$\beta$
  & MCA demand weight
  & 0.10 & 0.25 & 0.25 \\
$K$
  & KNN affinity count
  & 12 & 12 & 12 \\
$\tau$
  & Candidate threshold multiplier
  & 1.0 & 1.0 & 1.0 \\
$\bar{d}_{\text{nn}} \cdot \tau$ (m)
  & Resulting threshold
  & 322 & 484 & 1{,}315 \\
$P_{\text{total}}$
  & Total opening budget
  & 42 & 42 & 42 \\
$\lambda$
  & Uncovered-demand penalty
  & 5{,}000 & 5{,}000 & 5{,}000 \\
\midrule
\multicolumn{2}{l}{Grid-search combinations evaluated} & \multicolumn{3}{r}{81 per instance} \\
\multicolumn{2}{l}{Calibration criterion} & \multicolumn{3}{r}{Pareto-optimal (coverage, distance)} \\
\bottomrule
\end{tabular}

\end{table}

\begin{figure}[t]
  \centering
  \includegraphics[width=0.65\textwidth]{pareto_frontier_dmax_smin}
  \caption{Pareto frontier of coverage rate versus weighted mean assignment
           distance for the student-area grid search over $D_{\max}$
           and $S_{\min}$.
           The selected configuration (marked) lies on the efficient frontier.}
  \label{fig:pareto}
\end{figure}
